% Gerado de relatorio.md por gerar_relatorio.py; não editar este arquivo isoladamente.
\documentclass[11pt,a4paper]{article}
\usepackage[T1]{fontenc}
\usepackage[utf8]{inputenc}
\usepackage[brazil]{babel}
\usepackage{lmodern}
\usepackage[margin=23mm]{geometry}
\usepackage{microtype}
\usepackage[table]{xcolor}
\usepackage{longtable,array,booktabs}
\usepackage{xurl}
\usepackage[colorlinks=true,linkcolor=blue!45!black,urlcolor=blue!45!black]{hyperref}
\definecolor{pendente}{HTML}{B42318}
\newcommand{\pendencia}[1]{\textcolor{pendente}{[#1]}}
\hypersetup{pdftitle={Diário Emocional Inteligênte},pdfauthor={Equipe 007 - Licensa para Executar}}
\setlength{\parindent}{0pt}
\setlength{\parskip}{6pt}
\setlength{\emergencystretch}{3em}
\setlength{\tabcolsep}{5pt}
\renewcommand{\arraystretch}{1.18}
\begin{document}
\begin{titlepage}
\vspace*{25mm}
{\large EQUIPE 007\par}
{\large Licensa para Executar\par}
\vspace{12mm}
{\Huge\bfseries Diário Emocional Inteligênte\par}
\vspace{6mm}
{\LARGE Especificação de requisitos, conformidade ético-legal e arquitetura de referência\par}
\vspace{14mm}
{\large Versão 2.1 / 19 de novembro de 2026\par}
\vspace{8mm}
Dois agentes locais sob supervisão da psicóloga.\par
Agente 1: Extrator. Agente 2: Analista Longitudinal e Copiloto RAG.\par
\vfill
Autoria institucional: Equipe 007 - Licensa para Executar.\par
Challenge 2 - Bem-estar.\par
\pendencia{Validações técnicas abertas aparecem em vermelho e entre colchetes.}\par
Especificação do produto pretendido; não representa declaração de implementação.\par
\end{titlepage}
\tableofcontents
\newpage
\textbf{Equipe 007 - Licensa para Executar | Especificação de requisitos, conformidade ético-legal e arquitetura de referência | Versão 2.1 | 19/11/2026}

A Equipe 007 apresenta o Diário Emocional Inteligênte, uma aplicação de apoio à captação e organização de relatos entre sessões para profissionais de psicologia, com foco inicial em pacientes jovens-adultos. O produto combina registros do paciente, análise privada da psicóloga e prontuário oficial em camadas separadas. A identidade visual não constitui promessa de sigilo técnico ou de competência clínica.

Esta versão consolida a Especificação de Requisitos v2.1 e o documento “Diário Emocional Inteligente — Insights Clínicos e Evolução da Arquitetura de Agentes”, além das decisões de produto neles registradas. Ela descreve o comportamento pretendido e não declara funcionalidades implementadas nem resultados de avaliação já obtidos. Requisitos definidos pelas fontes são apresentados como referência; escolhas técnicas ainda não fechadas e critérios quantitativos a validar aparecem em vermelho e entre colchetes.

\section{Propósito, problema e objetivos}
O acompanhamento entre sessões chega ao consultório de forma fragmentada: datas, horários, acontecimentos fora da rotina e combinados podem se perder, e o paciente pode lembrar de situações importantes apenas na sessão seguinte. O produto propõe captar relatos voluntários de jovens-adultos, organizá-los em um workspace isolado por paciente e período e apoiar a preparação de pautas descritivas com acesso ao texto original.

O sistema facilita a consulta longitudinal, a organização de temas e subtemas, o acompanhamento de tarefas combinadas entre sessões e a recuperação de fontes normativas autorizadas. Saídas automatizadas são sugestões descritivas revisáveis; interpretação e decisão clínica permanecem exclusivamente com a psicóloga.

\begingroup\small
\begin{longtable}{>{\raggedright\arraybackslash}p{\dimexpr0.23\linewidth-2\tabcolsep\relax}>{\raggedright\arraybackslash}p{\dimexpr0.36\linewidth-2\tabcolsep\relax}>{\raggedright\arraybackslash}p{\dimexpr0.41\linewidth-2\tabcolsep\relax}}
\toprule
\textbf{Objetivo} & \textbf{Resultado esperado} & \textbf{Como verificar} \\ \midrule\endfirsthead
\toprule
\textbf{Objetivo} & \textbf{Resultado esperado} & \textbf{Como verificar} \\ \midrule\endhead
\bottomrule\endfoot
OBJ01 - Captar e localizar relatos & Reunir registros entre sessões em workspace isolado por paciente e período & Busca nos intervalos de 7, 14 e 30 dias, com acesso ao texto original \\[3pt]
OBJ02 - Estruturar registros & Organizar relatos em seis macrotemas, subtemas, modelo ABC e humor por cores & Esquema válido, trechos literais localizados e revisão humana preservada \\[3pt]
OBJ03 - Apoiar acompanhamento longitudinal & Comparar janelas temporais, acompanhar tarefas e produzir pauta descritiva & Cálculos determinísticos, coocorrências factuais e validação profissional \\[3pt]
OBJ04 - Proteger e rastrear informações & Separar camadas documentais e ancorar consultas em fontes autorizadas & Isolamento por camada, pseudonimização e citações verificáveis \\[3pt]
OBJ05 - Sustentar rigor acadêmico & Construir e avaliar classificador, RAG e dois agentes locais & Experimentos reproduzíveis e rastreabilidade MLOps \\[3pt]
\end{longtable}
\endgroup
A pesquisa de campo reportada nas fontes de requisitos reúne seis profissionais de psicologia (R1 a R6); seus achados orientam a concepção, mas não demonstram eficácia. A leitura profissional em menos de dois minutos é um objetivo de UX a validar. Não há alegação de melhora clínica, diagnóstico ou eficácia terapêutica do produto.

\section{Identidade, atores e fronteiras}
A autoria institucional é da \textbf{Equipe 007 - Licensa para Executar}. O produto se chama \textbf{Diário Emocional Inteligênte}. A aplicação destina-se à psicóloga e a pacientes jovens-adultos, com adesão voluntária e transparência sobre o tratamento dos dados.

\begingroup\small
\begin{longtable}{>{\raggedright\arraybackslash}p{\dimexpr0.23\linewidth-2\tabcolsep\relax}>{\raggedright\arraybackslash}p{\dimexpr0.36\linewidth-2\tabcolsep\relax}>{\raggedright\arraybackslash}p{\dimexpr0.41\linewidth-2\tabcolsep\relax}}
\toprule
\textbf{Ator} & \textbf{Responsabilidade no produto} & \textbf{Limite} \\ \midrule\endfirsthead
\toprule
\textbf{Ator} & \textbf{Responsabilidade no produto} & \textbf{Limite} \\ \midrule\endhead
\bottomrule\endfoot
Psicóloga & Consultar relatos, revisar sugestões, registrar notas privadas e validar evolução de prontuário & Única responsável pela interpretação clínica; autoria humana sempre identificada \\[3pt]
Paciente jovem-adulto & Registrar relatos voluntários, escolher o que compartilhar, informar humor e acompanhar tarefas & Não vê a análise privada nem temas/subtemas internos; pode acessar o prontuário oficial conforme as regras aplicáveis \\[3pt]
Agente 1 — Extrator & Estruturar cada relato com evidências literais & Não conversa com o paciente, não interpreta função de comportamento e não diagnostica \\[3pt]
Agente 2 — Analista longitudinal e copiloto RAG & Produzir sínteses factuais, rascunhos revisáveis e briefings normativos & Não estabelece causalidade, não diagnostica e não fala com o paciente \\[3pt]
Equipe 007 e orientação & Desenvolver, avaliar e validar requisitos acadêmicos & Não substituem revisão profissional nem análise jurídica \\[3pt]
\end{longtable}
\endgroup
O paciente controla voluntariamente quais registros compartilha com a psicóloga. O acesso a registros originais, análises privadas e prontuário oficial segue as três camadas descritas na seção 10; não há envio ou resumo automático para terceiros.

É vedada a emissão automatizada de diagnóstico, prognóstico, prescrição, recomendação terapêutica ou inferência clínica. O sistema não envia alertas externos autônomos. Materiais de TCC podem ser consultados apenas para apoio informativo autorizado; não autorizam o agente a formular condutas clínicas. As fontes ético-profissionais e legais de referência estão listadas na seção 21 (F14–F18).

\section{Escopo do produto e da entrega acadêmica}
O escopo consolidado reúne diário voluntário com texto livre (limite de 2.400 caracteres), campos TCC opcionais, seleção de humor, tarefas combinadas, pseudonimização local, seis macrotemas, extração ABC, análise longitudinal e consulta RAG normativa. As capacidades são requisitos-alvo; sua menção não afirma que estejam implementadas.

\begingroup\small
\begin{longtable}{>{\raggedright\arraybackslash}p{\dimexpr0.23\linewidth-2\tabcolsep\relax}>{\raggedright\arraybackslash}p{\dimexpr0.36\linewidth-2\tabcolsep\relax}>{\raggedright\arraybackslash}p{\dimexpr0.41\linewidth-2\tabcolsep\relax}}
\toprule
\textbf{Horizonte} & \textbf{Conteúdo} & \textbf{Evidência necessária} \\ \midrule\endfirsthead
\toprule
\textbf{Horizonte} & \textbf{Conteúdo} & \textbf{Evidência necessária} \\ \midrule\endhead
\bottomrule\endfoot
Especificação v2.1 & Requisitos, limites ético-legais, arquitetura de dois agentes locais e camadas documentais & Rastreabilidade entre fontes, requisitos e critérios de aceite \\[3pt]
MVP integrado & Fluxo local com pseudonimização, registro, extração, classificação, tarefas e consulta longitudinal/RAG & Testes ponta a ponta com dados fictícios, isolamento e falhas explícitas \\[3pt]
Evolução do produto & Interface acessível, leitura pré-consulta rápida, edição e prontuário validado & Avaliações técnicas, de usabilidade e revisão profissional \\[3pt]
\end{longtable}
\endgroup
As demonstrações acadêmicas devem utilizar dados fictícios. O uso de dados reais depende de controles de segurança, consentimento informado e revisão profissional e jurídica antes de qualquer piloto.

As restrições acadêmicas compreendem dois agentes colaborando, base de conhecimento, MLOps, rede neural construída e avaliada, API utilizada pelos agentes, versionamento, ética, explicabilidade e MVP funcional. A demonstração deve distinguir requisito, implementação e evidência obtida.

\section{Elicitação e organização das decisões}
Os requisitos incorporam a pesquisa qualitativa reportada com seis psicólogos e um parecer técnico-regulatório descrito nas fontes. Esses materiais orientam o desenho, mas não substituem validação contínua com profissionais, análise jurídica ou consentimento informado.

Os pontos ainda sujeitos a validação são a escolha exata de modelos e versões, dataset e licença, metas de avaliação, limites de segurança operacional e fluxos de consentimento e acesso. Exemplos de desenvolvimento devem ser fictícios.

Requisitos RF descrevem funções; RNF descrevem qualidades; RB descrevem regras de negócio; CA descrevem critérios de aceite. Alterações de taxonomia, modelo ou política de contagem devem gerar novas versões e indicar o impacto em resultados anteriores.

As escolhas consolidadas e as validações restantes estão resumidas na seção 20.

\section{Jornadas de uso}
\subsection{Abrir o contexto de uma pessoa}
A profissional pode começar no modo geral e solicitar a abertura de um paciente. O sistema resolve o cadastro por identificador e dados autorizados. Quando houver homônimos ou ambiguidade, apresenta opções antes de acessar os registros. Ao entrar no projeto, mostra explicitamente a pessoa selecionada e o período da consulta.

Dentro de um projeto, a profissional cria novas conversas ou retoma conversas anteriores. As mensagens pertencem à conversa; arquivos e registros pertencem ao projeto. Criar um novo chat não apaga nem duplica o histórico documental do paciente. Consultas gerais, como recebimentos recentes, podem reunir metadados dos pacientes autorizados sem misturar seus relatos em uma única síntese pessoal.

\subsection{Registrar um relato por texto ou voz}
A paciente registra voluntariamente um relato de até 2.400 caracteres, pode preencher campos estruturados opcionais inspirados em TCC, selecionar uma cor de humor e marcar tarefas combinadas. Entradas maiores são recusadas explicitamente, sem truncamento silencioso. O paciente decide quais entradas compartilhar.

Microperguntas opcionais podem ajudar a iniciar o relato sem induzir respostas nem exigir preenchimento de campos estruturados.

Para áudio, a transcrição ocorre localmente em memória e o arquivo de voz é descartado após o processamento. Texto original permanece preservado; a IA não o sobrescreve. A camada de processamento cria uma cópia pseudonimizada antes de encaminhar conteúdo aos modelos locais. Falhas e incertezas devem ser explícitas.

\subsection{Explorar temas e subtemas}
O dashboard apresenta assuntos identificados nos registros do período. A psicóloga pode navegar de tema para subtema e abrir as entradas que sustentam a organização. Subtemas gerados automaticamente ficam disponíveis para edição; correções, renomeações e uniões mantêm histórico. O paciente não visualiza essa classificação no produto.

\subsection{Consultar e agir pelo chat}
A psicóloga consulta registros em janelas de 7, 14 ou 30 dias, acompanha tarefas e solicita sínteses factuais ou briefings normativos. O Agente 2 aponta trechos e fontes, abstém-se quando a evidência é insuficiente e nunca converte coocorrência em causalidade. Rascunhos de evolução só passam ao prontuário oficial após revisão, edição e validação explícita da psicóloga.

\section{Requisitos funcionais}
\begingroup\small
\begin{longtable}{>{\raggedright\arraybackslash}p{\dimexpr0.23\linewidth-2\tabcolsep\relax}>{\raggedright\arraybackslash}p{\dimexpr0.36\linewidth-2\tabcolsep\relax}>{\raggedright\arraybackslash}p{\dimexpr0.41\linewidth-2\tabcolsep\relax}}
\toprule
\textbf{ID} & \textbf{Comportamento requerido} & \textbf{Critério de aceite associado} \\ \midrule\endfirsthead
\toprule
\textbf{ID} & \textbf{Comportamento requerido} & \textbf{Critério de aceite associado} \\ \midrule\endhead
\bottomrule\endfoot
RF01 & Manter workspace isolado por paciente, com UUID estável e anamnese restrita fora do pipeline de IA & CA01: homônimos não se misturam e dados de anamnese não chegam aos prompts \\[3pt]
RF02 & Receber relato de até 2.400 caracteres, campos TCC opcionais e cor de humor, aplicando pseudonimização local & CA02: excesso é recusado sem truncamento; Regex/NER substituem identificadores localmente \\[3pt]
RF03 & Sugerir seis macrotemas com MLP próprio e representar humor em dimensão visual separada & CA03: tema é sugestão revisável; humor não aparece como macrotema \\[3pt]
RF04 & Extrair situação, pensamento, sentimento, comportamento e consequência (ABC), de 1 a 6 subtemas, polaridade e indícios textuais de atenção com citação literal & CA04: toda citação corresponde ao original e toda saída inválida é rejeitada \\[3pt]
RF05 & Gerenciar status de tarefas combinadas e timeline por janelas temporais & CA05: estado e contagens são determinísticos nas janelas de 7, 14 e 30 dias \\[3pt]
RF06 & Produzir sínteses longitudinais, rascunhos de evolução e briefings RAG com o Agente 2 & CA06: sem causalidade/diagnóstico; fonte com score mínimo de 0,70 citada por documento e seção, ou abstenção \\[3pt]
RF07 & Isolar registro do paciente, análise privada da psicóloga e prontuário oficial em três camadas & CA07: análise e rascunhos não validados nunca aparecem no prontuário \\[3pt]
RF08 & Oferecer interface acessível, escaneável e com leitura profissional pré-consulta em menos de dois minutos & CA08: teclado, alto contraste e tipografia ampliada disponíveis; identidade visual personalizável \\[3pt]
\end{longtable}
\endgroup
Os critérios são requisitos de aceite, não resultados já obtidos. O desenho prevê adesão voluntária, seleção dos registros compartilhados e ausência de módulo de envio ou resumo automático para terceiros.

\section{Regras do produto e requisitos de qualidade}
\begingroup\small
\begin{longtable}{>{\raggedright\arraybackslash}p{\dimexpr0.23\linewidth-2\tabcolsep\relax}>{\raggedright\arraybackslash}p{\dimexpr0.36\linewidth-2\tabcolsep\relax}>{\raggedright\arraybackslash}p{\dimexpr0.41\linewidth-2\tabcolsep\relax}}
\toprule
\textbf{ID} & \textbf{Regra ou qualidade} & \textbf{Evidência de verificação} \\ \midrule\endfirsthead
\toprule
\textbf{ID} & \textbf{Regra ou qualidade} & \textbf{Evidência de verificação} \\ \midrule\endhead
\bottomrule\endfoot
RB01 & Vedação a diagnóstico, prognóstico, prescrição, recomendação terapêutica ou inferência clínica automatizada & Pedidos dessa natureza recebem delimitação de escopo e resposta descritiva \\[3pt]
RB02 & Registros do paciente, análise privada e prontuário oficial permanecem em três camadas sem compartilhamento externo automático & Testes de autorização confirmam isolamento e revisão humana antes de validação do prontuário \\[3pt]
RB03 & Relato original, anotação profissional e sugestão da IA têm autoria visualmente distinta & Metadados de autoria presentes; edição de tema não altera o original \\[3pt]
RB04 & Sinalização de atenção é passiva, apresenta estado e trecho literal e não aciona terceiros & Nenhuma comunicação automática; sinalização limitada a presente, indeterminado ou não identificado \\[3pt]
RB05 & Coocorrência temporal não significa causalidade clínica & Sínteses usam descrições factuais, sem concluir causa, evolução clínica ou diagnóstico \\[3pt]
RNF01 & Processamento dos dados clínicos local e offline, sem envio a APIs externas & Testes de rede e de prompt injection; nenhum tráfego clínico externo \\[3pt]
RNF02 & Rastreabilidade e reprodutibilidade & Versões, digest, latência e execução registrados sem texto bruto em logs \\[3pt]
RNF03 & Isolamento por paciente, camada e finalidade & Dados de outra pessoa e anamnese não chegam ao contexto dos agentes \\[3pt]
RNF04 & Integridade e recuperação & Escritas, backup e restauração verificáveis; políticas ainda exigem validação \\[3pt]
RNF05 & Acessibilidade e clareza & Teclado, alto contraste, fontes ampliadas e leitura que não depende apenas de cor \\[3pt]
RNF06 & Desempenho e qualidade observáveis & Latência, falhas, recursos e tarefas medidos com critérios publicados \\[3pt]
RNF07 & Reprodutibilidade experimental & Dados, licença, divisões, sementes, configurações e dependências identificados \\[3pt]
\end{longtable}
\endgroup
Processamento local reduz exposição a serviços externos, mas não comprova segurança ou conformidade por si só. Controles de autenticação, proteção de arquivos, backup, retenção e exclusão precisam ser implementados e verificados antes do uso com dados reais.

\section{Arquitetura de referência e responsabilidades}
O sistema possui dois percursos complementares. Na ingestão, recebe o original, transforma o conteúdo quando necessário, classifica, persiste e atualiza os índices. Na consulta, resolve o contexto, seleciona uma operação, recupera dados e apresenta uma resposta ou ação rastreável. O processamento de um diário não depende de um chat estar aberto.

\begingroup\small
\begin{longtable}{>{\raggedright\arraybackslash}p{\dimexpr0.23\linewidth-2\tabcolsep\relax}>{\raggedright\arraybackslash}p{\dimexpr0.36\linewidth-2\tabcolsep\relax}>{\raggedright\arraybackslash}p{\dimexpr0.41\linewidth-2\tabcolsep\relax}}
\toprule
\textbf{Componente lógico} & \textbf{Responsabilidade} & \textbf{Saída/controle} \\ \midrule\endfirsthead
\toprule
\textbf{Componente lógico} & \textbf{Responsabilidade} & \textbf{Saída/controle} \\ \midrule\endhead
\bottomrule\endfoot
Camada determinística & Ingestão, pseudonimização, validação, cálculos, persistência e visualização & Operações reproduzíveis em Python/SQL, sem depender de geração de linguagem \\[3pt]
Classificador temático & MLP próprio sobre embeddings SBERT & Sugestão de um macrotema, pontuação e versão \\[3pt]
Agente 1 — Extrator & Processar cada novo relato em esquema fechado & ABC, subtemas, sentimento/polaridade e marcador passivo com evidências \\[3pt]
Agente 2 — Analista longitudinal e copiloto RAG & Analisar janelas e responder consultas normativas autorizadas & Sínteses, rascunhos revisáveis ou briefing citado/abstenção \\[3pt]
Interface e camadas documentais & Controlar contexto, consentimento e visualização por perfil & Acesso separado a registro, análise privada e prontuário \\[3pt]
\end{longtable}
\endgroup
A arquitetura adota estritamente dois agentes locais. O Agente 1 é acionado por registro e opera com temperatura zero, semente fixa e esquema fechado; extrai situação, pensamento, sentimento, comportamento e consequência apenas quando sustentados pelo texto. O Agente 2 reúne análise longitudinal e RAG, sem criar um terceiro agente para cada função. A colaboração ocorre por dados tipados e evidências validadas.

Operações críticas de privacidade, cálculo, validação e escrita são determinísticas. Os agentes executam localmente por Ollama; nenhuma API externa recebe dados clínicos. Modelos candidatos incluem Qwen 2.5, Qwen3 4B e Gemma 3 4B; a versão final e seus requisitos de hardware ainda precisam de decisão. Temperatura zero e semente fixa são parâmetros de execução, não garantia universal de determinismo.

\section{Ferramentas e controle das ações}
A interface expõe aos agentes apenas operações locais delimitadas, com esquemas fechados, validação de argumentos e identidade de paciente verificada. Nenhum agente tem acesso irrestrito ao computador, à rede externa ou a operações de compartilhamento.

\begingroup\small
\begin{longtable}{>{\raggedright\arraybackslash}p{\dimexpr0.23\linewidth-2\tabcolsep\relax}>{\raggedright\arraybackslash}p{\dimexpr0.36\linewidth-2\tabcolsep\relax}>{\raggedright\arraybackslash}p{\dimexpr0.41\linewidth-2\tabcolsep\relax}}
\toprule
\textbf{Família} & \textbf{Exemplos} & \textbf{Controle} \\ \midrule\endfirsthead
\toprule
\textbf{Família} & \textbf{Exemplos} & \textbf{Controle} \\ \midrule\endhead
\bottomrule\endfoot
Privacidade & Detectar identificadores, pseudonimizar e resolver tokens autorizados & Execução local; tabela de associação cifrada e fora dos prompts \\[3pt]
Extração & Validar esquema ABC, subtemas, polaridade e marcadores & Pydantic e validação exata das citações no original \\[3pt]
Histórico & Filtrar janelas, contar entradas, cores e tarefas & Python/SQL determinístico com paciente e período explícitos \\[3pt]
RAG & Recuperar normas e técnicas autorizadas & Busca híbrida, referência de documento/seção e abstenção \\[3pt]
Prontuário & Preparar rascunho de evolução & Migração apenas após revisão e validação explícita da psicóloga \\[3pt]
\end{longtable}
\endgroup
O texto do paciente e documentos recuperados são dados, nunca instruções capazes de alterar permissões ou o contrato dos agentes. Pedidos diagnósticos e marcadores de atenção não iniciam ações externas nem substituem avaliação humana.

\section{Projetos, arquivos e modelo de dados}
Na entrada, regras locais detectam CPF, RG, telefones, e-mails, endereços e dados bancários e os substituem por marcadores neutros. NER local pseudonimiza nomes, locais e instituições com tokens consistentes por paciente. A proposta cita `validate-docbr`, Microsoft Presidio com spaCy (`pt\_core\_news\_lg`) ou BERTimbau NER; são alternativas técnicas a validar, não componentes já implantados.

Os registros são vinculados a um UUID estável por paciente. A anamnese e a tabela que relaciona identidade real a pseudônimos ficam fora do pipeline de IA; a tabela de associação deve ser cifrada, com AES-256 como referência e gestão de chaves separada, e acessível somente no ambiente autenticado da psicóloga. A arquitetura separa dados por finalidade e camada:

\begingroup\small
\begin{longtable}{>{\raggedright\arraybackslash}p{\dimexpr0.23\linewidth-2\tabcolsep\relax}>{\raggedright\arraybackslash}p{\dimexpr0.36\linewidth-2\tabcolsep\relax}>{\raggedright\arraybackslash}p{\dimexpr0.41\linewidth-2\tabcolsep\relax}}
\toprule
\textbf{Camada} & \textbf{Conteúdo} & \textbf{Acesso e regra} \\ \midrule\endfirsthead
\toprule
\textbf{Camada} & \textbf{Conteúdo} & \textbf{Acesso e regra} \\ \midrule\endhead
\bottomrule\endfoot
1. Registro do paciente & Texto original, campos TCC opcionais, cor de humor e tarefas & Paciente e psicóloga conforme compartilhamento; original nunca é sobrescrito pela IA \\[3pt]
2. Análise da psicóloga & ABC extraído, temas, subtemas, marcadores, sínteses, briefings RAG e notas privadas & Acesso exclusivo da psicóloga; saída automática é sugestão editável \\[3pt]
3. Prontuário oficial & Registro documental e rascunhos de evolução revisados & Paciente tem acesso conforme regra aplicável; nada migra sem edição e validação explícita da psicóloga \\[3pt]
\end{longtable}
\endgroup
Não há quarta camada de compartilhamento externo nem resumo automático para terceiros. A estrutura física, banco e política de backup ainda devem ser definidos sem comprometer esse isolamento.

\begingroup\small
\begin{longtable}{>{\raggedright\arraybackslash}p{\dimexpr0.23\linewidth-2\tabcolsep\relax}>{\raggedright\arraybackslash}p{\dimexpr0.36\linewidth-2\tabcolsep\relax}>{\raggedright\arraybackslash}p{\dimexpr0.41\linewidth-2\tabcolsep\relax}}
\toprule
\textbf{Entidade} & \textbf{Conteúdo mínimo} & \textbf{Relações} \\ \midrule\endfirsthead
\toprule
\textbf{Entidade} & \textbf{Conteúdo mínimo} & \textbf{Relações} \\ \midrule\endhead
\bottomrule\endfoot
Profissional & ID e configurações autorizadas & Projetos e biblioteca geral \\[3pt]
Projeto de paciente & ID, identificação e estado & Conversas, registros e anexos \\[3pt]
Conversa & ID, escopo geral ou projeto, mensagens & Fontes consultadas e execuções \\[3pt]
Registro & ID, autor/origem, texto e datas & Anexos, classificações e versões \\[3pt]
Arquivo & ID, caminho controlado, tipo, hash e origem & Páginas, áudio ou extração de texto \\[3pt]
Classificação & Registro, taxonomia, versão de modelo e saída & Temas e eventuais correções \\[3pt]
Subtema & ID, projeto, tema pai, nome e estado & Entradas de apoio e histórico de edição \\[3pt]
Trecho recuperável & Fonte, versão, localização e conteúdo & Índice lexical/vetorial \\[3pt]
Execução de ferramenta & ID, contexto, operação, resultado e estado & Alterações e fontes utilizadas \\[3pt]
\end{longtable}
\endgroup
Índices, resumos e gráficos são derivados reconstruíveis; não substituem o original. Alterações devem manter histórico e invalidar derivados incompatíveis. A tabela de associação entre pseudônimos e identidades deve permanecer cifrada e separada dos dados enviados aos agentes.

A trilha registra versões de modelo e taxonomia, digest e latência, sem guardar texto bruto do paciente em logs. A política de retenção, backup e recuperação ainda precisa ser formalizada antes do uso com dados reais.

\section{Rede neural para classificação temática}
A rede classifica o assunto do relato em uma categoria principal entre seis macrotemas. Humor, sentimento e polaridade são dimensões distintas, não classes temáticas. A psicóloga pode revisar a sugestão.

\begingroup\small
\begin{longtable}{>{\raggedright\arraybackslash}p{\dimexpr0.5\linewidth-2\tabcolsep\relax}>{\raggedright\arraybackslash}p{\dimexpr0.5\linewidth-2\tabcolsep\relax}}
\toprule
\textbf{Macrotema} & \textbf{Escopo descritivo} \\ \midrule\endfirsthead
\toprule
\textbf{Macrotema} & \textbf{Escopo descritivo} \\ \midrule\endhead
\bottomrule\endfoot
Autoestima e Identidade & Autopercepção, inseguranças e posicionamento pessoal \\[3pt]
Relacionamentos e Família & Vínculos amorosos, amizades e dinâmicas familiares \\[3pt]
Trabalho, Estudos e Finanças & Pressões e realizações profissionais, acadêmicas e materiais \\[3pt]
Saúde, Sono e Rotina & Condição física, sono, medicações e eventos fora da rotina \\[3pt]
Perdas, Luto e Rupturas & Falecimentos, términos, demissões e mudanças bruscas \\[3pt]
Outros / Indeterminado & Classe residual quando nenhuma categoria temática supera o limiar definido \\[3pt]
\end{longtable}
\endgroup
O MLP usa embeddings SBERT como representação e sugere uma categoria quando a confiança é de pelo menos 0,50; abaixo disso, retorna Outros / Indeterminado. Humor é atributo separado, selecionado pelo paciente e apresentado em escala visual com opções como calma, alegria, tristeza, medo, preocupação, frustração e raiva; não deve ser inferido como macrotema. A classificação de uma categoria principal segue o enquadramento multiclasse, não multilabel (F09).

O requisito acadêmico é construir e treinar o MLP do projeto, documentando camadas, inicialização, passagem direta, perda, retropropagação e atualização de pesos. O uso de SBERT como representação pré-treinada não substitui a implementação e avaliação da rede classificadora. \pendencia{Escolher o modelo SBERT, dimensões, configuração do MLP, otimizador, taxa de aprendizado, batch e regra de parada.}

Com V atributos de entrada, H unidades ocultas e C categorias, uma rede densa de uma camada oculta possui V vezes H + H + H vezes C + C parâmetros treináveis. O vocabulário deve ser ajustado somente no treino. Pesos são aprendidos; hiperparâmetros são escolhidos na validação. A implementação deve incluir verificação numérica de gradientes ou outra evidência de correção da retropropagação antes de interpretar resultados semânticos. Referência didática: F10.

\section{Dataset, anotação e experimentos da rede}
O conjunto de treinamento precisa conter textos e referências compatíveis com a taxonomia. Tamanho, idioma e presença de uma coluna de categoria não bastam para selecionar uma base. Devem ser avaliados domínio, licença, possibilidade de uso, distribuição das classes, ambiguidade e proximidade entre exemplos.

\begingroup\small
\begin{longtable}{>{\raggedright\arraybackslash}p{\dimexpr0.23\linewidth-2\tabcolsep\relax}>{\raggedright\arraybackslash}p{\dimexpr0.36\linewidth-2\tabcolsep\relax}>{\raggedright\arraybackslash}p{\dimexpr0.41\linewidth-2\tabcolsep\relax}}
\toprule
\textbf{Campo proposto} & \textbf{Finalidade} & \textbf{Uso no aprendizado} \\ \midrule\endfirsthead
\toprule
\textbf{Campo proposto} & \textbf{Finalidade} & \textbf{Uso no aprendizado} \\ \midrule\endhead
\bottomrule\endfoot
id & Identificar o exemplo & Metadado \\[3pt]
texto & Conteúdo do relato ou trecho & Entrada, após representação \\[3pt]
tema(s) & Categoria(s) esperada(s) & Alvo supervisionado \\[3pt]
grupo\_origem & Pessoa, entrevista ou família de templates & Separação dos conjuntos; não é pista de classificação \\[3pt]
origem\_rotulo & Humano, regra ou modelo, com versão & Qualidade e interpretação das referências \\[3pt]
split & Treino, validação ou teste & Controle do experimento \\[3pt]
origem/licença & Proveniência e condições de uso & Auditoria de dados \\[3pt]
\end{longtable}
\endgroup
\pendencia{Selecionar dataset, fonte, versão, licença, quantidade por classe e uso de dados sintéticos ou públicos autorizados.} O conjunto deve corresponder aos seis macrotemas, excluir identificadores e documentar origem e licença. O MuPe pode ser explorado como corpus narrativo, mas não contém automaticamente a taxonomia do produto (F07); o MASSIVE trata de intenções de assistentes e não é substituto direto de relatos temáticos (F08). Qualquer adaptação exige justificativa e controle de vazamento entre pessoas ou templates.

Definir anotação, revisão independente e tratamento de desacordo. Rótulos gerados por LLM ou regras são supervisão fraca: métricas sobre eles medem concordância com a referência, não validade clínica nem desempenho em uso real.

Separar treino, validação e teste antes de ajustar representação ou modelos. Variações quase idênticas e relatos do mesmo grupo devem ser controlados para evitar vazamento. As proporções, critérios de agrupamento, deduplicação e tamanho mínimo de teste por classe permanecem por definir; a seleção da configuração usa validação e o teste final fica reservado.

Comparar a rede com uma referência simples, como a classe majoritária, e com um baseline textual compatível. Apresentar curvas de treino e validação, matriz de confusão, precisão/recall por macrotema e macro-F1. Definir limiares de aceite, tratamento de classes raras, repetições e reporte de incerteza antes de interpretar resultados.

Baixo desempenho também no treino sugere investigar capacidade, representação, otimização e qualidade dos rótulos. Desempenho alto no treino e inferior na validação sugere investigar generalização. Não ajustar tudo simultaneamente: registrar experimentos com alterações controladas e conservar resultados negativos.

\section{Subtemas pessoais e análise longitudinal}
Subtemas organizam assuntos específicos dentro dos seis macrotemas e não alteram a taxonomia do classificador. O Agente 1 propõe de um a seis subtemas por relato, sempre ancorados em trecho literal e posições inicial/final verificáveis.

O Agente 2 compara registros da mesma pessoa em janelas de 7, 14 e 30 dias. Cálculos de contagem, distribuição de humor/polaridade, timeline e tarefas são executados em Python/SQL. Sínteses podem apontar coocorrências e sequência temporal em linguagem factual (por exemplo, “ocorreu antes de” ou “coincidiu com”), sem afirmar causalidade, melhora, piora ou intensidade clínica.

O status das tarefas é determinístico (feita/não feita) e aparece na timeline. Toda síntese informa janela e conjunto de registros. Variação de contagens pode refletir frequência de envio ou mudança de versão do classificador, não uma mudança clínica.

\section{RAG, memória e conhecimento profissional}
RAG combina recuperação de informações externas com geração de resposta. A memória persistente é o conjunto de registros; o RAG seleciona partes pertinentes desse conjunto para uma consulta. Não é necessário retreinar o LLM quando um arquivo é adicionado ao índice. A formulação original é referência técnica histórica, não um resource recente da entrega (F06).

\begingroup\small
\begin{longtable}{>{\raggedright\arraybackslash}p{\dimexpr0.23\linewidth-2\tabcolsep\relax}>{\raggedright\arraybackslash}p{\dimexpr0.36\linewidth-2\tabcolsep\relax}>{\raggedright\arraybackslash}p{\dimexpr0.41\linewidth-2\tabcolsep\relax}}
\toprule
\textbf{Coleção} & \textbf{Conteúdo} & \textbf{Regra de uso} \\ \midrule\endfirsthead
\toprule
\textbf{Coleção} & \textbf{Conteúdo} & \textbf{Regra de uso} \\ \midrule\endhead
\bottomrule\endfoot
Instruções da aplicação & Papel, limites e ferramentas permitidas & Aplicadas diretamente; não dependem de busca por similaridade \\[3pt]
Espaço geral profissional & Materiais, preferências e metodologias autorizadas & Contextualização da profissional, sem virar fato sobre pacientes \\[3pt]
Registros por paciente & Diários, notas, anexos e transcrições & Recuperação com identidade e autorização obrigatórias \\[3pt]
Histórico de conversa & Mensagens do chat e contexto recente & Continuidade; não equivale a registro profissional confirmado \\[3pt]
Biblioteca publicada & Documentos externos com autoria e condições de uso & Citação e distinção entre evidência publicada e relato individual \\[3pt]
\end{longtable}
\endgroup
Na ingestão: preservar o arquivo; extrair texto; identificar documento, página ou intervalo; dividir em trechos quando necessário; indexar; verificar disponibilidade. Na consulta: resolver escopo; aplicar filtros; recuperar candidatos; selecionar contexto; gerar resposta; verificar referências e apresentar limitações.

O corpus normativo/técnico autorizado inclui Código de Ética Profissional, Resolução CFP nº 001/2009, notas CFP/CRP-SP de 2025 e técnicas aprovadas. Uma norma ou publicação contextualiza o trabalho profissional, mas não é evidência sobre um paciente. Fonte incompleta, não autorizada ou não pertinente não deve ser citada como suporte.

O desenho de referência usa busca híbrida BM25 + embeddings e5-small, fusão por Reciprocal Rank Fusion (RRF) e limiar de score de pelo menos 0,70. O multilingual-e5-small é uma opção documentada, sujeita à validação no corpus (F11). O briefing cita documento e seção; abaixo do limiar ou sem fonte sustentadora, o agente declara abstenção. Os parâmetros precisam ser avaliados no corpus autorizado antes de serem tratados como desempenho comprovado.

Filtros de autorização e escopo são aplicados antes da recuperação. O classificador temático não pode excluir silenciosamente registros relevantes. Um erro de recuperação deve resultar em fonte insuficiente ou abstenção, não em conteúdo inventado.

\section{Tempo, visualização e experiência profissional}
O dashboard e a análise profissional devem ser escaneáveis em menos de dois minutos, com fontes ampliadas e fluxo simplificado para acessibilidade, inclusive para pessoas com menor familiaridade digital. As janelas de referência são 7, 14 e 30 dias. Pessoa, período, número de registros e origem automática/revisada permanecem visíveis.

O humor selecionado pelo paciente é mostrado por cores como dimensão separada dos temas, com legenda textual acessível. Uma cor não deve representar gravidade nem intensidade clínica. A interface prevê navegação por teclado, alto contraste e identidade visual personalizável. As métricas devem permitir abrir as entradas de origem.

\section{Evidências de engenharia agêntica e aplicação ao projeto}
A pesquisa bibliográfica deste relatório considera fontes até \textbf{19/11/2026}. Os trabalhos citados fundamentam critérios de projeto; não certificam este produto nem constituem um padrão único do MIT. Recomendações de engenharia de fornecedores são identificadas como tal, sem pressupor adoção de seus serviços. Requisitos clínico-regulatórios e de UX também incorporam os dois documentos de produto listados na seção 21.

\begingroup\small
\begin{longtable}{>{\raggedright\arraybackslash}p{\dimexpr0.23\linewidth-2\tabcolsep\relax}>{\raggedright\arraybackslash}p{\dimexpr0.36\linewidth-2\tabcolsep\relax}>{\raggedright\arraybackslash}p{\dimexpr0.41\linewidth-2\tabcolsep\relax}}
\toprule
\textbf{Fonte} & \textbf{Achado relevante} & \textbf{Aplicação ao produto} \\ \midrule\endfirsthead
\toprule
\textbf{Fonte} & \textbf{Achado relevante} & \textbf{Aplicação ao produto} \\ \midrule\endhead
\bottomrule\endfoot
F01 - AI Agent Index, MIT, artigo de 2026 & Identifica lacunas de transparência na documentação de agentes implantados & Documentar ferramentas, autonomia, limites, modelos e resultados de testes \\[3pt]
F02 - ImpactBench, 2026 & Propõe avaliação de comportamentos ligados a impacto humano, usando simulações, juiz LLM e revisão especializada & Avaliar autonomia da profissional, clareza de limites e possibilidade de conferir fontes; não presumir benefício clínico \\[3pt]
F03 - Building effective agents, 2024 & Distingue fluxos previsíveis de decisões dinâmicas por modelos e recomenda complexidade justificada & Usar código para cálculos e persistência; justificar cada papel agêntico \\[3pt]
F04 - Writing effective tools, 2025 & Enfatiza ferramentas bem descritas, retornos úteis e avaliação & Contratos tipados, escopo claro, erros explícitos e retornos suficientes \\[3pt]
F05 - Demystifying evals, 2026 & Distingue resposta do agente, trajetória de ações e estado final & Verificar o que foi realmente gravado, lido ou alterado; repetir cenários \\[3pt]
\end{longtable}
\endgroup
O AI Agent Index descreve sistemas e sua transparência; não fornece certificação de segurança para uma implementação particular. O ImpactBench é um preprint e infraestrutura de avaliação em evolução. Seu uso de revisão especializada impede interpretar a presença de um juiz LLM como prova de que revisão humana é dispensável. A aplicação desses trabalhos ao produto é uma decisão de engenharia da Equipe 007.

A avaliação deve cobrir tanto componentes quanto o fluxo integrado. Um modelo que classifica corretamente pode estar conectado à pessoa errada; uma resposta bem escrita pode citar uma fonte que não a sustenta; uma mensagem de sucesso pode ocultar falha de gravação. Por isso, o teste precisa observar dados, operações e resultado final.

\section{Plano de avaliação e critérios de aceite}
\begingroup\small
\begin{longtable}{>{\raggedright\arraybackslash}p{\dimexpr0.23\linewidth-2\tabcolsep\relax}>{\raggedright\arraybackslash}p{\dimexpr0.36\linewidth-2\tabcolsep\relax}>{\raggedright\arraybackslash}p{\dimexpr0.41\linewidth-2\tabcolsep\relax}}
\toprule
\textbf{Camada} & \textbf{Evidência de avaliação} & \textbf{Condição de reporte} \\ \midrule\endfirsthead
\toprule
\textbf{Camada} & \textbf{Evidência de avaliação} & \textbf{Condição de reporte} \\ \midrule\endhead
\bottomrule\endfoot
Rede temática & Macro-F1, precisão/recall por classe, matriz e curvas & Taxonomia, divisão de dados e origem dos rótulos declaradas \\[3pt]
Subtemas & Aderência de evidências, duplicação, edição e estabilidade & Sem afirmar descoberta de condição clínica \\[3pt]
Recuperação & Recall@k, HitRate@k e casos sem fonte relevante & Recall e acerto de ao menos um resultado separados \\[3pt]
Resposta & Fonte correta, suporte factual e reconhecimento de lacuna & ID de fonte existente não basta: trecho precisa sustentar a afirmação \\[3pt]
Ferramentas & Argumentos, autorização, cancelamento e estado final & Não avaliar apenas a frase de confirmação \\[3pt]
Usabilidade & Conclusão de tarefas, tempo e correções necessárias & Perfil, amostra e roteiro documentados \\[3pt]
Operação & Latência, erros, recursos e recuperação & Máquina, versão e limites da medição declarados \\[3pt]
\end{longtable}
\endgroup
As métricas, rubricas, critérios de aceite e número de repetições precisam ser definidos antes do teste final. Resultados sem referência independente devem ser identificados como preliminares; concordância entre modelos não estabelece verdade.

\begingroup\small
\begin{longtable}{>{\raggedright\arraybackslash}p{\dimexpr0.23\linewidth-2\tabcolsep\relax}>{\raggedright\arraybackslash}p{\dimexpr0.36\linewidth-2\tabcolsep\relax}>{\raggedright\arraybackslash}p{\dimexpr0.41\linewidth-2\tabcolsep\relax}}
\toprule
\textbf{Cenário de aceite} & \textbf{Preparação sintética} & \textbf{Resultado esperado} \\ \midrule\endfirsthead
\toprule
\textbf{Cenário de aceite} & \textbf{Preparação sintética} & \textbf{Resultado esperado} \\ \midrule\endhead
\bottomrule\endfoot
CA01 - Isolamento & Dois perfis fictícios com nomes iguais e anamneses & Nenhuma mistura; anamnese não chega ao agente \\[3pt]
CA02 - Limite e pseudonimização & Relatos de 2.400 e 2.401 caracteres com PII sintética & Limite respeitado sem truncamento; identificadores pseudonimizados localmente \\[3pt]
CA03 - Tema e humor & Um relato com tema e humor informado & Tema revisável e humor apresentado em dimensão distinta \\[3pt]
CA04 - Citação literal & Citação válida e citação inventada em saída de teste & A primeira é localizada; a segunda invalida a saída \\[3pt]
CA05 - Tarefas e períodos & Tarefas feitas/não feitas em janelas conhecidas & Timeline e contagens coincidem com cálculo determinístico \\[3pt]
CA06 - RAG insuficiente & Pergunta sem fonte acima do limiar & Agente se abstém, sem inventar referência ou resposta \\[3pt]
CA07 - Camadas & Rascunho não validado e análise privada & Nenhum conteúdo migra ou aparece no prontuário oficial \\[3pt]
CA08 - Acessibilidade & Navegação por teclado, alto contraste e texto ampliado & Funções essenciais permanecem utilizáveis e legíveis \\[3pt]
RB04 - Marcador passivo & Trecho sintético contém marcador de atenção & Estado e trecho aparecem apenas no painel; nenhum contato externo \\[3pt]
RB05 - Coocorrência & Eventos em semanas sucessivas & Texto descreve sequência/coocorrência sem atribuir causalidade \\[3pt]
\end{longtable}
\endgroup
A demonstração inicial utiliza pelo menos dois perfis fictícios para testar isolamento. Uma execução bem-sucedida isolada não é suficiente para declarar confiabilidade do sistema.

\section{MLOps, versionamento e evolução}
Cada experimento deve registrar commit de código, versão/hash dos dados, taxonomia, forma de divisão, configuração, modelo, métricas e artefatos. A reprodução deve distinguir ambiente, dependências, semente e limitações de determinismo do hardware e runtime.

Mudanças no contrato afetam classificador, persistência, subtemas, RAG, dashboard e testes. A taxonomia deve possuir versão, e resultados históricos devem indicar qual versão foi utilizada. Reprocessamento precisa ser uma operação explícita e auditável.

Documentos e índices também têm ciclo de vida: ingestão, validação, disponibilidade, atualização e remoção. Monitorar registros que não foram indexados e índices desatualizados. O chat deve comunicar quando uma resposta usa um conjunto ainda incompleto.

\section{Rastreabilidade acadêmica, riscos e validação}
\begingroup\small
\begin{longtable}{>{\raggedright\arraybackslash}p{\dimexpr0.23\linewidth-2\tabcolsep\relax}>{\raggedright\arraybackslash}p{\dimexpr0.36\linewidth-2\tabcolsep\relax}>{\raggedright\arraybackslash}p{\dimexpr0.41\linewidth-2\tabcolsep\relax}}
\toprule
\textbf{Restrição/necessidade} & \textbf{Requisitos relacionados} & \textbf{Evidência esperada} \\ \midrule\endfirsthead
\toprule
\textbf{Restrição/necessidade} & \textbf{Requisitos relacionados} & \textbf{Evidência esperada} \\ \midrule\endhead
\bottomrule\endfoot
Dois agentes colaborando & RF04, RF06 & Papéis distintos, execução local e troca de dados estruturados \\[3pt]
Base de conhecimento & RF06 & Corpus autorizado, ingestão rastreável, recuperação híbrida e citação \\[3pt]
Rede neural própria & RF03 & MLP próprio, operações de treino e avaliação por macrotema \\[3pt]
MLOps & RNF02, RNF07; seção 18 & Execução rastreável e artefatos reproduzíveis \\[3pt]
API por agentes & RF04, RF06 & Integração local com tratamento de falhas e esquema validado \\[3pt]
Ética, privacidade e explicabilidade & RB01 a RB05, RNF01 a RNF03 & Limites, camadas, origem, pseudonimização e citações verificáveis \\[3pt]
MVP e versionamento & RF01 a RF08, RNF01 a RNF07 & Fluxo integrado com dados fictícios, testes e histórico de alterações \\[3pt]
\end{longtable}
\endgroup
Referências acadêmicas e normativas apoiam a especificação, mas não certificam a implementação. Dados reais não devem ser usados antes de validar consentimento, segurança, regras profissionais e requisitos legais pertinentes.

\begingroup\small
\begin{longtable}{>{\raggedright\arraybackslash}p{\dimexpr0.23\linewidth-2\tabcolsep\relax}>{\raggedright\arraybackslash}p{\dimexpr0.36\linewidth-2\tabcolsep\relax}>{\raggedright\arraybackslash}p{\dimexpr0.41\linewidth-2\tabcolsep\relax}}
\toprule
\textbf{Risco} & \textbf{Consequência} & \textbf{Tratamento} \\ \midrule\endfirsthead
\toprule
\textbf{Risco} & \textbf{Consequência} & \textbf{Tratamento} \\ \midrule\endhead
\bottomrule\endfoot
Vazamento de identidade & Exposição de dados sensíveis & Pseudonimização local, UUID por paciente, cofre cifrado separado e testes de isolamento \\[3pt]
Citação inventada ou fora de contexto & Síntese enganosa & Validar trechos literais contra o original e RAG contra documento/seção \\[3pt]
Marcador de atenção interpretado como diagnóstico & Falso alarme ou decisão clínica indevida & Exibição passiva com trecho literal; nenhuma ação ou comunicação automática \\[3pt]
Coocorrência interpretada como causalidade & Inferência clínica não sustentada & Linguagem factual, períodos explícitos e revisão da psicóloga \\[3pt]
Falha de agente ou modelo & Registro inconsistente ou saída inválida & Esquema fechado, validação, falha explícita e preservação do original \\[3pt]
Mistura entre análise e prontuário & Acesso indevido ou registro não validado & Três camadas isoladas e validação humana antes da migração \\[3pt]
\end{longtable}
\endgroup
A versão documenta decisões de produto já consolidadas e mantém separadas as escolhas técnicas e de avaliação ainda abertas.

\section{Decisões consolidadas e validações pendentes}
\subsection{Definições adotadas na especificação v2.1}
\hangindent=1em - Público inicial: profissionais de psicologia e pacientes jovens-adultos; adesão voluntária e transparência sobre o uso dos dados.

\hangindent=1em - Entrada: texto de até 2.400 caracteres, campos TCC opcionais, seleção de humor e status de tarefas; áudio transcrito localmente em memória e descartado após a transcrição.

\hangindent=1em - Privacidade: identificadores diretos removidos por regras locais; nomes, locais e instituições pseudonimizados consistentemente; UUID e tabela de associação cifrada mantida separada.

\hangindent=1em - Organização: seis macrotemas revisáveis, humor visual independente, extração ABC literal e subtemas com evidências localizáveis.

\hangindent=1em - Agentes: dois agentes locais — Agente 1 Extrator e Agente 2 Analista Longitudinal \& Copiloto RAG — apoiados por cálculos determinísticos.

\hangindent=1em - Acesso: registro do paciente, análise privada da psicóloga e prontuário oficial em três camadas; sem compartilhamento ou resumo automático a terceiros.

\hangindent=1em - Limites: nenhuma saída diagnóstica, prescritiva, prognóstica, causal ou terapêutica; marcador textual passivo sem alerta externo.

\hangindent=1em - Consulta: janelas de 7, 14 e 30 dias; RAG híbrido BM25 + e5-small com RRF, citação documento/seção e abstenção abaixo do score especificado.

\subsection{Validações técnicas e operacionais ainda necessárias}
\pendencia{Selecionar dataset, licença, processo de anotação e divisão de grupos; estabelecer baseline, métricas e critérios de aceite do classificador.}

\pendencia{Definir versão exata do SBERT e dos LLMs locais, arquitetura e hiperparâmetros do MLP, runtime, hardware mínimo e plano de indisponibilidade.}

\pendencia{Validar tamanho de trechos, implementação e calibração do limiar RAG, cobertura do corpus e testes de abstenção.}

\pendencia{Formalizar autenticação, consentimento, permissões para compartilhar cada registro, retenção, backup, restauração e exclusão antes de qualquer piloto com dados reais.}

\pendencia{Definir rubricas, perfis e tarefas de usabilidade, repetições, tolerância a falhas e metas mensuráveis de qualidade e desempenho.}

\pendencia{Registrar responsáveis, aprovação de requisitos e evidências produzidas para cada versão do MVP.}

\pendencia{Validar acompanhamento longitudinal sem diagnóstico ou tratamento automático entre sessões.}

\pendencia{Avaliar mudanças temporais e comorbidades sem confundir diagnósticos ou alterar condutas clínicas.}

\pendencia{Construir ou selecionar dados culturalmente nativos em português e auditar modelos abertos quanto a licença, proveniência e viés.}

\pendencia{Testar diversidade de transtornos e abordagens terapêuticas, incluindo referências a DBT e terapia humanista, sem prescrição automatizada.}

\section{Referências e critérios de uso}
Data de corte desta versão: 19/11/2026. As fontes a seguir fundamentam requisitos, limites ético-profissionais e decisões de engenharia. A inclusão bibliográfica não significa adoção de ferramenta, licença para redistribuir seu conteúdo ou validação do produto.

F01. STAUFER, Leon et al. The 2025 AI Agent Index: Documenting Technical and Safety Features of Deployed Agentic AI Systems. Preprint submetido em 19/02/2026; versão 2 de 06/05/2026. Portal hospedado no MIT. \url{https://arxiv.org/abs/2602.17753v2} . Portal: \url{https://aiagentindex.mit.edu/} . Contribuição: documentação de autonomia, ferramentas e avaliação; não prescreve uma arquitetura única.

F02. ARCHIWARANGUPROK, Chayapatr et al. Open Benchmark of AI Impact on Humans (ImpactBench): an expert-guided, open platform for the holistic evaluation of AI impact on humans. Preprint, versão 1, 14/10/2026. \url{https://arxiv.org/abs/2610.20690v1} . Plataforma MIT Media Lab: \url{https://impactbench.media.mit.edu/about} . Contribuição: avaliação de comportamentos e impacto humano; inclui limitações de simulação e julgamento automatizado, sem validar este aplicativo.

F03. ANTHROPIC. Building effective agents. 19/12/2024. \url{https://www.anthropic.com/research/building-effective-agents} . Contribuição: fluxos, agentes e simplicidade arquitetural. É orientação de engenharia de fornecedor; detalhes de ferramentas evoluem.

F04. ANTHROPIC. Writing effective tools for AI agents. 11/09/2025. \url{https://www.anthropic.com/research/writing-tools-for-agents} . Contribuição: descrição, limites, retorno e avaliação das ferramentas.

F05. ANTHROPIC. Demystifying evals for AI agents. 09/01/2026. \url{https://www.anthropic.com/research/demystifying-evals-for-ai-agents} . Contribuição: avaliação da execução e do estado final, além do texto produzido.

F06. LEWIS, Patrick et al. Retrieval-Augmented Generation for Knowledge-Intensive NLP Tasks. NeurIPS, 2020. \url{https://proceedings.neurips.cc/paper/2020/hash/6b493230205f780e1bc26945df7481e5-Abstract.html} . Referência histórica do método, fora da janela de dois anos para resource recente.

F07. LEAL, Sidney Evaldo et al. MuPe Life Stories Dataset: Spontaneous Speech in Brazilian Portuguese with a Case Study Evaluation on ASR Bias against Speakers Groups and Topic Modeling. COLING, 2025. \url{https://aclanthology.org/2025.coling-main.407/} . Contribuição: candidato de domínio narrativo; uso e adaptação ainda dependem de decisão.

F08. AMAZON SCIENCE. MASSIVE, documentação do dataset. \url{https://huggingface.co/datasets/AmazonScience/massive} . Contribuição: exemplo de classificação de intenções; rótulos não equivalem à taxonomia temática do produto.

F09. SCIKIT-LEARN. Multiclass and multioutput algorithms; multilabel classification. \url{https://scikit-learn.org/stable/modules/multiclass.html} . Contribuição: distinção entre categorias exclusivas e simultâneas e representação das referências.

F10. NUMPY COMMUNITY. Deep learning on MNIST. \url{https://numpy.org/numpy-tutorials/tutorial-deep-learning-on-mnist/} . Contribuição: exemplo didático de rede construída com operações NumPy; dados de imagens não substituem o dataset textual do projeto.

F11. INTFLOAT. multilingual-e5-small, model card. \url{https://huggingface.co/intfloat/multilingual-e5-small} . Contribuição: candidato de embeddings, limites e preparação das entradas; seleção depende de benchmark próprio.

F12. EQUIPE 007 - LICENSA PARA EXECUTAR. Especificação de Requisitos — Diário Emocional Inteligente, versão 2.1 consolidada, 19/11/2026. \url{https://docs.google.com/document/d/1Vd2MqRL89ZywdTZhE9LsI2x8bMcIKYzP/edit} . Fonte-base dos requisitos de produto, público, camadas e critérios de aceite.

F13. EQUIPE 007 - LICENSA PARA EXECUTAR. Diário Emocional Inteligente — Insights Clínicos e Evolução da Arquitetura de Agentes, 2026. \url{https://docs.google.com/document/d/1RXIFI4TdtOksuJmNa12xX5T9e4pUaN4esNH7i80Vsng/edit} . Fonte de recomendações sobre pseudonimização, agentes, taxonomia, supervisão e arquitetura documental.

F14. BRASIL. Lei nº 13.709, de 14 de agosto de 2018. Lei Geral de Proteção de Dados Pessoais (LGPD), arts. 5º, 6º, 13 e 46.

F15. CONSELHO FEDERAL DE PSICOLOGIA. Resolução CFP nº 010/2005. Código de Ética Profissional do Psicólogo, arts. 9º e 10º.

F16. CONSELHO FEDERAL DE PSICOLOGIA. Resolução CFP nº 001/2009. Registro documental decorrente da prestação de serviços psicológicos, incluindo art. 6º.

F17. CONSELHO FEDERAL DE PSICOLOGIA. Nota de Posicionamento sobre o Uso de Inteligência Artificial na Prática Psicológica, 2025.

F18. CONSELHO REGIONAL DE PSICOLOGIA DE SÃO PAULO. Nota Orientativa sobre o Uso de Ferramentas de Inteligência Artificial na Psicologia, 2025.

F19. NA, Hongbin; HUA, Yining; WANG, Zimu; SHEN, Tao; YU, Beibei; WANG, Lilin; WANG, Wei; TOROUS, John; CHEN, Ling. A Survey of Large Language Models in Psychotherapy: Current Landscape and Future Directions. Findings of ACL, 2025. \url{https://arxiv.org/pdf/2502.11095v1} . Contribuição: destaca integração fragmentada, sintomas evolutivos e comorbidades, vieses linguísticos/culturais, concentração em transtornos comuns e necessidade de diversidade terapêutica. Os desafios são agenda futura; não comprovam eficácia deste aplicativo.

Esta versão alinha requisitos e arquitetura às fontes F12 e F13. Implementação, conformidade e eficácia precisam ser demonstradas por avaliações próprias; a especificação não substitui parecer profissional ou jurídico.

\end{document}
